\chapter{Otto 計算と勾配流}
\label{ch:seminar-09-otto}

第~\ref{ch:seminar-08-benamou-brenier}章で，$\Wass_2^2$ が
「連続の方程式に従う流れの最小作用」であることを見た．
この事実は，確率測度の空間 $\Pp_2(\R^d)$ を
\textbf{無限次元の（形式的な）リーマン多様体}とみなす視点を示唆する：
各点 $\rho$ における「動く方向」は流れの初速度であり，
その「長さ」は運動エネルギーで測られ，
$\Wass_2$ はこの計量が定める測地距離にほかならない．
この形式的幾何を整備したのが Otto の計算である．

本章の目標は二つある．第一に，$\Pp_2$ の\textbf{接空間}と\textbf{Otto 計量}を
定め，曲線の速度・勾配を測度の空間で語れるようにすること．
第二に，この計量のもとで\textbf{エネルギー汎関数の勾配流}を考えると，
相対エントロピーの勾配流が\textbf{熱方程式}に一致する，という
Jordan--Kinderlehrer--Otto の結果に到達することである．
これは確率測度の動力学を変分原理で捉える枠組みであり，
章末では拡散モデルとの接続にも触れる．
なお本章は形式的なリーマン幾何として議論し，
厳密化（存在・収束）は文献に譲る．

%% ============================================================
%% 9.1 接空間と Otto 計量
%% ============================================================
\section{接空間と Otto 計量}
\label{sec:sem-otto-metric}

$\Pp_2(\R^d)$ 内の曲線 $(\rho_t)$ の「速度」を定める．
連続の方程式（定義~\ref{def:sem-continuity-equation}）
$\partial_t\rho_t+\diverg(\rho_t v_t)=0$ は，
分布の変化 $\partial_t\rho_t$ を速度場 $v_t$ に結びつける．
ただし同じ $\partial_t\rho_t$ を与える $v_t$ は一意でない
（$\diverg(\rho_t w)=0$ なる $w$ を加えてもよい）．
運動エネルギー $\int\norm{v_t}^2\d\rho_t$ を最小にする代表を選ぶと，
それは\textbf{勾配場} $v_t=\nabla\psi_t$ になる．これが接ベクトルの正準な表示を与える．

\begin{claim}{接空間（Helmholtz 分解）}{sem-tangent-space}
  $\rho\in\Pp_2(\R^d)$ における接空間を
  \[
    \Tan_\rho\Pp_2
    \defeq \overline{\bigl\{\,\nabla\psi : \psi\in\Cc_c^\infty(\R^d)\,\bigr\}}^{\,L^2(\rho)}
  \]
  （$L^2(\rho;\R^d)$ における勾配場の閉包）と定める．
  任意のベクトル場 $v\in L^2(\rho;\R^d)$ は
  \[
    v = \nabla\psi + w,
    \qquad \nabla\psi\in\Tan_\rho\Pp_2,\quad \diverg(\rho w)=0,
  \]
  と直交分解され（$\rho$-重み付き Helmholtz 分解），
  $\nabla\psi$ は連続の方程式の制約
  $\diverg(\rho v)=\diverg(\rho\nabla\psi)$ を保ったまま
  $\int\norm{\,\cdot\,}^2\d\rho$ を最小にする代表である．
\end{claim}

\begin{definition}{Otto 計量}{sem-otto-inner}
  曲線の速度 $\partial_t\rho$ を，連続の方程式を通じて
  接ベクトル $\nabla\psi\in\Tan_\rho\Pp_2$ と同一視する：
  $\partial_t\rho=-\diverg(\rho\nabla\psi)$．
  二つの接ベクトル $\nabla\psi_1,\nabla\psi_2$ に対する\textbf{Otto 計量}を
  \[
    \inner{\nabla\psi_1}{\nabla\psi_2}_\rho
    \defeq
    \int_{\R^d}\inner{\nabla\psi_1(x)}{\nabla\psi_2(x)}\,\d\rho(x)
  \]
  で定める．対応するノルムは $\norm{\nabla\psi}_\rho^2=\int\norm{\nabla\psi}^2\d\rho$．
\end{definition}

\begin{remark}{$\Wass_2$ は Otto 計量の測地距離}{sem-otto-w2}
  曲線 $(\rho_t)$ の長さを $\int_0^1\norm{\partial_t\rho_t}_{\rho_t}\,\d t$ で測ると，
  これは第~\ref{ch:seminar-08-benamou-brenier}章の作用 $\mathcal{A}(\rho,v)$
  （最小代表 $v=\nabla\psi$ をとったもの）の平方根の時間積分にあたり，
  Benamou--Brenier の定理~\ref{thm:sem-benamou-brenier}は
  \[
    \Wass_2(\alpha,\beta)
    = \inf\Bigl\{\textstyle\int_0^1\norm{\partial_t\rho_t}_{\rho_t}\,\d t
      : \rho_0=\alpha,\ \rho_1=\beta\Bigr\}
  \]
  と読み替えられる．すなわち $\Wass_2$ は Otto 計量が定める\textbf{測地距離}であり，
  測地線は前章の変位補間である．これが「$\Pp_2$ をリーマン多様体とみなす」
  という形式的視点の正当化（の骨子）である．
\end{remark}

%% ============================================================
%% 9.2 第一変分と Wasserstein 勾配
%% ============================================================
\section{第一変分と Wasserstein 勾配}
\label{sec:sem-w2-gradient}

リーマン多様体上の勾配流を語るには，汎関数の\textbf{勾配}が要る．
まず汎関数の微分（第一変分）を定め，それを Otto 計量で「勾配」に翻訳する．

\begin{definition}{第一変分}{sem-first-variation}
  汎関数 $F:\Pp_2(\R^d)\to\R$ の $\rho$ における\textbf{第一変分}
  $\dfrac{\delta F}{\delta\rho}:\R^d\to\R$ とは，
  質量 $0$ の摂動 $\chi$（$\int\chi\,\d x=0$）に対して
  \[
    \left.\frac{\d}{\d s}\right|_{s=0} F(\rho+s\chi)
    = \int_{\R^d} \frac{\delta F}{\delta\rho}(x)\,\chi(x)\,\d x
  \]
  を満たす関数をいう（存在するとき）．
\end{definition}

ユークリッド空間では勾配は微分そのものだが，
リーマン多様体では計量を通して接空間の元へ「持ち上げる」必要がある．
Otto 計量のもとでの勾配は次で与えられる．

\begin{claim}{Wasserstein 勾配}{sem-w2-grad-formula}
  $F$ の Otto 計量に関する\textbf{勾配} $\mathrm{grad}_{\Wass_2}F(\rho)$ は，
  接ベクトル $\nabla\bigl(\tfrac{\delta F}{\delta\rho}\bigr)$ に対応する：
  曲線方向 $\partial_t\rho=-\diverg(\rho\nabla\psi)$ に対して
  \[
    \left.\frac{\d}{\d t}\right|_{0}F(\rho_t)
    = \inner{\,\mathrm{grad}_{\Wass_2}F(\rho)\,}{\nabla\psi}_\rho,
    \qquad
    \mathrm{grad}_{\Wass_2}F(\rho)=\nabla\Bigl(\tfrac{\delta F}{\delta\rho}\Bigr).
  \]
\end{claim}

\begin{proof}[形式的な確認]
  連続の方程式と部分積分，および第一変分の定義より
  \[
    \frac{\d}{\d t}F(\rho_t)
    = \int \frac{\delta F}{\delta\rho}\,\partial_t\rho_t
    = -\int \frac{\delta F}{\delta\rho}\,\diverg(\rho_t\nabla\psi)
    = \int \inner{\nabla\Bigl(\tfrac{\delta F}{\delta\rho}\Bigr)}{\nabla\psi}\,\d\rho_t
    = \inner{\nabla\Bigl(\tfrac{\delta F}{\delta\rho}\Bigr)}{\nabla\psi}_{\rho_t}.
  \]
  これは勾配の定義（方向微分が計量内積で書ける）にほかならない．
\end{proof}

%% ============================================================
%% 9.3 JKO スキームと勾配流
%% ============================================================
\section{JKO スキームと勾配流}
\label{sec:sem-jko}

ユークリッド空間の勾配流 $\dot x=-\nabla F(x)$ は，
\textbf{陰的 Euler 法}（最小化運動）
$x^{k+1}=\argmin_x\bigl(\tfrac{1}{2\tau}\norm{x-x^k}^2+F(x)\bigr)$ の
$\tau\to0$ 極限として特徴づけられる．
距離 $\norm{x-x^k}$ を $\Wass_2$ に置き換えると，測度の空間での勾配流が定義できる．
これが Jordan--Kinderlehrer--Otto（JKO）の最小化運動スキームである．

\begin{definition}{JKO スキーム}{sem-jko-scheme}
  汎関数 $F:\Pp_2(\R^d)\to\R$，時間刻み $\tau>0$，初期測度 $\rho^0$ に対し，
  \[
    \rho^{k+1}
    \in
    \argmin_{\rho\in\Pp_2(\R^d)}
    \left\{\,
      \frac{1}{2\tau}\,\Wass_2^2(\rho,\rho^k) + F(\rho)
    \,\right\}
    \qquad (k=0,1,2,\dots)
  \]
  により列 $(\rho^k)$ を定める．各ステップは
  「$\rho^k$ から $\Wass_2$ で近く，かつ $F$ を下げる」測度を選ぶ操作である．
\end{definition}

\begin{claim}{JKO の収束と勾配流}{sem-jko-limit}
  $F$ が適当な条件（下半連続・変位凸性など）を満たすとき，
  時間刻み $\tau\to0$ で，JKO 列の区分定数補間
  $\rho^\tau(t)$ は曲線 $(\rho_t)$ に収束し，その極限は
  $F$ の\textbf{Wasserstein 勾配流}
  \[
    \partial_t\rho_t
    = -\,\mathrm{grad}_{\Wass_2}F(\rho_t)
    = \diverg\!\Bigl(\rho_t\,\nabla\tfrac{\delta F}{\delta\rho}(\rho_t)\Bigr)
  \]
  を満たす（主張~\ref{clm:sem-w2-grad-formula}を代入した形）．
\end{claim}

最後の偏微分方程式の右辺は，連続の方程式
$\partial_t\rho+\diverg(\rho v)=0$ において
速度を $v=-\nabla\tfrac{\delta F}{\delta\rho}$ にとったものである：
勾配流とは「自由エネルギー $F$ を最も急に下げる向きに質量を流す」ことにほかならない．

%% ============================================================
%% 9.4 熱方程式はエントロピーの勾配流
%% ============================================================
\section{熱方程式はエントロピーの勾配流}
\label{sec:sem-heat-as-gradient-flow}

最も基本的な汎関数として\textbf{（負の）エントロピー}
（Lebesgue 測度に対する相対エントロピー）
\[
  \Hm(\rho) \defeq \int_{\R^d} \rho(x)\log\rho(x)\,\d x
\]
をとる．その Wasserstein 勾配流が熱方程式になることを示す．

\begin{theorem}{Jordan--Kinderlehrer--Otto}{sem-jko-heat}
  $F=\Hm$ の Wasserstein 勾配流は\textbf{熱方程式}である：
  \[
    \partial_t\rho_t = \Delta\rho_t .
  \]
  したがって熱拡散は，エントロピーを $\Wass_2$ 計量のもとで
  最も急に減少させる流れとして特徴づけられる．
\end{theorem}

\begin{proof}[形式的な計算]
  $\Hm$ の第一変分を計算する．$\frac{\d}{\d s}\int(\rho+s\chi)\log(\rho+s\chi)
  \big|_{s=0}=\int(\log\rho+1)\chi$ なので，
  $\tfrac{\delta\Hm}{\delta\rho}=\log\rho+1$．よって
  \[
    \nabla\tfrac{\delta\Hm}{\delta\rho}
    = \nabla\log\rho = \frac{\nabla\rho}{\rho}.
  \]
  主張~\ref{clm:sem-jko-limit}の勾配流に代入すると
  \[
    \partial_t\rho
    = \diverg\!\Bigl(\rho\,\nabla\log\rho\Bigr)
    = \diverg\!\Bigl(\rho\cdot\frac{\nabla\rho}{\rho}\Bigr)
    = \diverg(\nabla\rho)
    = \Delta\rho.
  \]
\end{proof}

\begin{remark}{ポテンシャル項と Fokker--Planck}{sem-fokker-planck}
  外場ポテンシャル $V:\R^d\to\R$ を加えた自由エネルギー
  $F(\rho)=\Hm(\rho)+\int V\,\d\rho$ をとると，
  $\tfrac{\delta F}{\delta\rho}=\log\rho+1+V$ より勾配流は
  \[
    \partial_t\rho
    = \diverg\!\bigl(\rho\,\nabla(\log\rho+V)\bigr)
    = \Delta\rho + \diverg(\rho\,\nabla V),
  \]
  すなわち\textbf{Fokker--Planck 方程式}となる．
  その定常解はギブス測度 $\rho_\infty\propto e^{-V}$ で，
  これは $F$ の最小点でもある（エントロピーと位置エネルギーの釣り合い）．
  確率の言葉では，これは過減衰 Langevin 拡散
  $\d X_t=-\nabla V(X_t)\,\d t+\sqrt2\,\d B_t$ の密度の発展である．
\end{remark}

%% ============================================================
%% 9.5 応用：拡散モデルと Schrödinger ブリッジ（おまけ）
%% ============================================================
\section{応用：拡散モデルと Schrödinger ブリッジ}
\label{sec:sem-otto-diffusion}

本節は本筋から離れ，勾配流の視点が\textbf{拡散モデル}と
\textbf{Schrödinger ブリッジ}にどう結びつくかを述べる補足である．
証明は与えず，対応のみを示す．

\begin{memo*}
  \textbf{スコアベース拡散モデルとの対応．}
  拡散モデルは，データ $\rho_0$ にノイズを徐々に加える\textbf{前進過程}
  （Fokker--Planck／Langevin，備考~\ref{rem:sem-fokker-planck}）と，
  それを時間反転してノイズからデータを生成する\textbf{後退過程}からなる．
  後退過程の駆動項にはスコア $\nabla\log\rho_t$ が現れるが，
  これはまさにエントロピー $\Hm$ の Wasserstein 勾配
  $\nabla\tfrac{\delta\Hm}{\delta\rho}=\nabla\log\rho_t$
  （定理~\ref{thm:sem-jko-heat}）である．
  すなわち拡散モデルの学習対象「スコア」は，
  \textbf{エントロピー勾配流の速度場}に等しい．
  Flow Matching（前章 §\ref{sec:sem-bb-flow-matching}）が
  最小作用の流れ（測地線）を学ぶのに対し，
  拡散モデルはエントロピー勾配流に沿う速度場を学ぶ，と整理できる．
\end{memo*}

\begin{memo*}
  \textbf{Schrödinger ブリッジとの対応．}
  第~\ref{ch:seminar-03-entropic}章・第~\ref{ch:seminar-04-sinkhorn}章の
  エントロピー正則化 OT は，動的に見ると
  「基準となるブラウン運動からの相対エントロピーを最小にしつつ
  端点 $\alpha,\beta$ をつなぐ確率過程を求める」
  \textbf{Schrödinger 問題}に対応する．
  これは Benamou--Brenier の作用に拡散項を加えた形（運動エネルギー＋エントロピー）であり，
  正則化パラメータ $\varepsilon\to0$ で動的最適輸送（前章）へ収束する．
  近年の\textbf{拡散 Schrödinger ブリッジ}は，この問題を
  反復射影（Sinkhorn の動的版）で解いて生成モデルに用いる枠組みであり，
  本セミナーのエントロピー正則化と動的定式化がともに部品として効く．
\end{memo*}

要するに，Otto 計算が与える「測度の動力学＝自由エネルギーの勾配流」という描像は，
測地線（Flow Matching）と拡散（スコア・Schrödinger）を
同じ変分原理のもとに並べる共通言語となる．
次章では，この勾配流の\textbf{凸性}——汎関数が測地線に沿って凸かどうか——を問い，
それが空間の\textbf{曲率}と結びつくことを見る．
